\chapter{Introduction}
\label{chapter:introducao}

Cyber-physical systems (CPS) connect software, networks, and physical processes in environments where performance and reliability must coexist. In critical infrastructures, a security failure can affect computational systems, physical assets, and service continuity \cite{Cardenas2009,Mitchell2014,Alcaraz2019}. Intrusion detection systems (IDS) must therefore handle high-dimensional data, time constraints, and operational traceability.

Feature selection reduces dimensionality and redundancy before classification \cite{Hall1999,Saeys2007,guyon2003introduction}. Original GRASP combines greedy randomized construction and local search \cite{ResendeRibeiro2010}. Its GRASP-FS adaptation for feature selection is addressed by \citeonline{quincozes2020grasp,quincozes2021performance,quincozes2022extended}. In the monolithic implementations studied here, ranking, construction, refinement, and execution records remain in one process, restricting separation of responsibilities and stage inspection.

This dissertation proposes \textit{G-FShield}, a distributed, event-driven architecture for GRASP-FS in \acp{ids} applied to \acp{cps}. \ac{drg} computes four independent rankings and constructs initial solutions; \ac{dls} combines VND or RVND with Bit-Flip, IWSS, and IWSSR to refine them. Apache Kafka connects the stages, and a web layer queries persisted events, metrics, states, identifiers, timestamps, and intermediate results. Observability here means inspecting the recorded external state of search, not causally explaining classifier decisions.

The general objective is to design, implement, and evaluate this distributed decomposition under a traceable protocol. The specific objectives are:

\begin{itemize}
  \item separate candidate generation, neighborhood control, local search, verification, and monitoring;
  \item provide a web layer for inspecting runs and persisted artifacts;
  \item compare quality, dimensionality reduction, time, and resources with monoliths and a baseline under common data, evaluator, and computational ceiling;
  \item preserve enough provenance to reconstruct runs and explicitly delimit properties not yet measured.
\end{itemize}

Four research questions guide the evaluation. \textbf{RQ1}: under the common protocol, what predictive quality and dimensionality reduction are obtained by G-FShield configurations, the monoliths, and the all-features baseline? \textbf{RQ2}: how are RCL ranking, neighborhood controller, and local-search preference descriptively related to final solutions and bounded search time? \textbf{RQ3}: what CPU and memory are observed for each architecture under the same aggregate ceiling, without conflating consumption with efficiency? \textbf{RQ4}: which events, identifiers, metrics, and artifacts allow a run to be reconstructed, and which observability and recovery properties remain unevaluated?

The contribution includes the implemented architecture, a formal 27-arm campaign, and a second causal campaign with 30 pairs that keeps RF--VND--IWSSR identical while varying distributed or sequential execution. The analysis separates quality, latency, throughput, parallelism, and resource cost instead of presuming architectural superiority. Both use a stratified training/validation/test split, a common final evaluator, and verifiable manifests. Corpus independence is architectural rather than experimental: input must be labeled tabular data compatible with the service contract. The evaluation covers only the ERENO corpus and J48 classifier.

The remainder of this dissertation is organized as follows. Chapter~\ref{chapter:fundamentacao} presents the foundations; Chapter~\ref{chapter:relacionados} discusses related work; Chapter~\ref{chapter:desenvolvimento} describes architecture and implementation; Chapter~\ref{chapter:avaliacao} details protocol and results; and Chapter~\ref{chapter:conclusao} summarizes conclusions, limitations, and future work.
